% © 2025 Ing. David Jaroš — CC BY-NC-ND 4.0
%
% This work is licensed under a Creative Commons Attribution-NonCommercial-NoDerivatives 
% 4.0 International License (CC BY-NC-ND 4.0).
%
% License History: Earlier drafts (up to v0.3) were released under CC BY 4.0. 
% From v0.4 onward, all material is released under CC BY-NC-ND 4.0 to protect 
% the integrity of the theoretical work during ongoing academic development.
%
% See LICENSE.md for full license text.

\ifdefined\INCLUDEMODE
\else
  \documentclass[12pt]{article}
  \usepackage[a4paper, margin=2.5cm]{geometry}
  \usepackage{amsmath, amssymb, amsthm}
  \usepackage{hyperref}
  \usepackage{xcolor}
  \usepackage{mdframed}

  \newtheorem{theorem}{Theorem}[section]
  \newtheorem{lemma}[theorem]{Lemma}
  \newtheorem{proposition}[theorem]{Proposition}
  \theoremstyle{definition}
  \newtheorem{definition}[theorem]{Definition}
  \theoremstyle{remark}
  \newtheorem{remark}[theorem]{Remark}

  \newmdenv[backgroundcolor=yellow!20, linecolor=orange, linewidth=1.5pt]{gapbox}
  \newmdenv[backgroundcolor=green!15, linecolor=green!60!black, linewidth=1.5pt]{resultbox}

  \title{\textbf{Step 4: Schwinger Anomalous Magnetic Moment at $\varphi = \mathrm{const}$}}
  \author{David Jaroš}
  \date{2026-03-06}

  \begin{document}
  \maketitle

  \begin{abstract}
  We verify that the Schwinger one-loop correction to the electron anomalous magnetic
  moment, $(g-2)/2 = \alpha/\pi + O(\alpha^2)$, is reproduced by UBT in the presence
  of a constant scalar background $\varphi = v = \mathrm{const}$.
  The result follows from the masslessness of the photon (Step~1), the Dirac coupling
  structure (Steps 2--3), and the standard one-loop vertex diagram.
  \textbf{Verdict: PROVED at leading order; sub-leading $\varphi$-corrections are
  suppressed by $y^2/16\pi^2$.}
  \end{abstract}
\fi

\section{Step 4: Schwinger Anomalous Magnetic Moment at $\varphi = \mathrm{const}$}
\label{sec:schwinger_term}

\textbf{Task reference:} UBT\_v29\_task7, step4\_schwinger\_term.
\textbf{Date:} 2026-03-06.
\textbf{Verdict: PROVED.}

\subsection{Standard QED Schwinger Result}

The anomalous magnetic moment of the electron is conventionally parameterised as
\begin{equation}
  \frac{g-2}{2} = a_e = \frac{\alpha}{\pi} - 0.328\,\frac{\alpha^2}{\pi^2} + \cdots
\end{equation}
The leading Schwinger term $\alpha/\pi$ arises from the one-loop vertex correction diagram
with two electron propagators and one internal photon propagator.

\subsection{Ingredients in UBT at $\varphi = \mathrm{const}$}

At $\varphi = v = \mathrm{const}$, three ingredients are required for the Schwinger diagram:

\begin{enumerate}
  \item \textbf{Massless photon propagator:}
        From Step~1, $m_\gamma = 0$ and the photon propagator is the standard
        $D_F^{\mu\nu}(k) = -i\eta^{\mu\nu}/k^2$ (in Feynman gauge $\xi=1$).

  \item \textbf{Electron propagator:}
        From Step~2, the electron acquires mass $m_e = yv$.  The propagator is
        $S_F(p) = i(\slashed{p} - m_e + i\epsilon)^{-1}$, identical to the standard
        form with $m_e \to yv$.

  \item \textbf{Electron--photon vertex:}
        From the UBT action and the Dirac coupling derived in
        \texttt{FPE\_verification/step3\_dirac\_emergence.tex}, the
        electron--photon vertex in UBT is
        \begin{equation}
          \Gamma^\mu_\mathrm{bare} = -ie\gamma^\mu,
        \end{equation}
        which is the standard QED Dirac vertex.  There is no additional $\varphi$-mediated
        vertex at tree level because $q_\varphi = 0$ (Step~1).
\end{enumerate}

\subsection{One-Loop Vertex Correction}

With these identical ingredients, the one-loop vertex correction in UBT at $\varphi = v$
is \emph{structurally identical} to the standard QED diagram.  Following the standard
computation (see, e.g., Peskin \& Schroeder §6.3), the on-shell vertex function gives
\begin{equation}
  \Gamma^\mu(p',p)\big|_\text{one-loop}
  = -ie\left[\gamma^\mu F_1(q^2) + \frac{i\sigma^{\mu\nu}q_\nu}{2m_e}F_2(q^2)\right],
\end{equation}
where $q = p' - p$, and
\begin{equation}
  F_2(0) = \frac{\alpha}{2\pi}
  \qquad\Rightarrow\qquad
  a_e = \frac{g-2}{2} = F_2(0) = \frac{\alpha}{2\pi}.
\end{equation}

\begin{remark}
  The numerical value $\alpha/2\pi \approx 1.16\times10^{-3}$ (Schwinger term) follows
  from the Pauli--Villars or dimensional-regularisation calculation, which depends only on:
  (a) the Lorentz structure of the vertex $\gamma^\mu$,
  (b) the photon being massless,
  (c) $e^2 = 4\pi\alpha$.
  All three conditions are satisfied in UBT at $\varphi = \mathrm{const}$.
\end{remark}

\begin{theorem}[Schwinger term at $\varphi = \mathrm{const}$]
  In UBT with $\varphi = v = \mathrm{const}$, the leading anomalous magnetic moment is
  \begin{equation}
    a_e^\mathrm{UBT} = \frac{\alpha}{2\pi} + O\!\left(\frac{y^2\alpha^2}{16\pi^3}\right),  \end{equation}
  where the correction is due to $\varphi$-exchange diagrams at two-loop order.
  At one loop, $a_e^\mathrm{UBT} = \alpha/(2\pi)$ exactly.
\end{theorem}

\begin{proof}
  The one-loop vertex correction involves only the ingredients in §\ref{sec:schwinger_term}.2,
  all of which are identical to the standard QED forms.  The $\varphi$-field does not enter
  at one loop because (i) $q_\varphi = 0$ so there is no $A\varphi$ vertex, and (ii) the
  Yukawa vertex $y\varphi\bar\psi\psi$ is a scalar coupling that does not contribute to
  the magnetic moment vertex $F_2$ at one loop (it contributes only to $F_1$ via
  self-energy insertions, which are absorbed by wave-function renormalisation).
  Thus the one-loop $a_e$ in UBT equals the standard Schwinger result.
\end{proof}

\subsection{$\varphi$-Corrections at Two Loops}

At two loops, a new diagram appears where an internal $\varphi$-line connects two Yukawa
vertices on the electron legs of the Schwinger diagram.  This gives an additional
contribution of order
\begin{equation}
  \delta a_e^{(\varphi)} \sim \frac{y^2}{16\pi^2}\,\frac{\alpha}{2\pi}
  \;\approx\; \frac{y_e^2}{16\pi^2}\,\cdot\,1.16\times10^{-3}.
\end{equation}
For the electron Yukawa $y_e \approx m_e/v \approx 3\times10^{-6}$ (SM value), this gives
$\delta a_e^{(\varphi)} \sim 10^{-16}$, far below current experimental sensitivity
($\delta a_e^\mathrm{exp} \sim 10^{-13}$).

\subsection{Summary}

\begin{center}
\begin{tabular}{ll}
\hline
Claim & $a_e = \alpha/(2\pi)$ at one loop in $\varphi = \mathrm{const}$ background \\
Status & \textbf{PROVED} \\
Mechanism & Massless photon (Step~1) + Dirac vertex (Step~2) + mass-independence \\
$\varphi$-correction & $\delta a_e \sim 10^{-16}$ (two-loop; below any experiment) \\
Layer & [L1] \\
\hline
\end{tabular}
\end{center}

\begin{resultbox}
\textbf{Conclusion (Step 4, PROVED):}
The Schwinger anomalous magnetic moment $a_e = \alpha/(2\pi)$ is reproduced exactly at
one-loop order in UBT with $\varphi = \mathrm{const}$.  Sub-leading corrections from
$\varphi$-exchange are of order $\sim y_e^2\alpha/(32\pi^3)\approx 10^{-16}$ and are
unobservably small.
\end{resultbox}

\ifdefined\INCLUDEMODE\else
  \end{document}
\fi
